\section*{Capítulo \ref{ch:g5:food-webs} --- Teias alimentares}

\begin{solution}{exo:g5:food-webs:1}
A rede de todas as \omterm{def:g3:food-chains:chain}{cadeias alimentares} de um \omterm{def:g2:where-animals-live:habitat}{hábitat} desenhadas ao
mesmo tempo --- cada \omterm{def:g5:biodiversity:species}{espécie} uma vez, uma seta para cada relação de
refeição. As cadeias se cruzam porque quase todos os \omterm{def:g1:animals-around-us:animal}{animais} comem
várias \omterm{def:g5:biodiversity:species}{espécies} e são comidos por várias.
\end{solution}

\begin{solution}{exo:g5:food-webs:2}
\omterm{def:g5:food-webs:roles}{Produtores}: as \omterm{def:g1:plants-around-us:plant}{plantas} (capim, trevo, \omterm{def:g1:plants-around-us:parts}{folhas} de carvalho).
\omterm{def:g5:food-webs:roles}{Consumidores}: os \omterm{def:g1:animals-around-us:animal}{animais} --- \omterm{def:g2:what-animals-eat:kinds}{herbívoros} como o coelho, \omterm{def:g2:what-animals-eat:kinds}{carnívoros}
como o gavião. \omterm{def:g5:food-webs:roles}{Decompositores}: os que apodrecem os restos mortos
(nenhum aparece na figura --- eles trabalham na sombra, embaixo de
tudo).
\end{solution}

\begin{solution}{exo:g5:food-webs:3}
O chapim come \omterm{ex:g2:animals-grow-up:butterfly}{lagartas} e gafanhotos; ele é comido pelo gavião-miúdo.
\end{solution}

\begin{solution}{exo:g5:food-webs:4}
Capim $\rightarrow$ gafanhoto $\rightarrow$ lagarto $\rightarrow$
gavião, e capim $\rightarrow$ gafanhoto $\rightarrow$ chapim
$\rightarrow$ gavião-miúdo.
\end{solution}

\begin{solution}{exo:g5:food-webs:5}
Muitos \omterm{def:g5:food-webs:roles}{produtores}, menos \omterm{def:g2:what-animals-eat:kinds}{herbívoros}, poucos \omterm{def:g2:what-animals-eat:kinds}{carnívoros}. Porque cada
andar se alimenta do de baixo: quem come nunca pode ser mais numeroso
do que aquilo que há para comer.
\end{solution}

\begin{solution}{exo:g5:food-webs:6}
Porque um casal de gaviões precisa de um campo inteiro de gafanhotos,
com vários lagartos no meio, para viver --- o topo da pirâmide se
apoia nos andares largos de baixo.
\end{solution}

\begin{solution}{exo:g5:food-webs:7}
Os caçadores deles --- a raposa, o gavião --- comem melhor e se
multiplicam em seguida; essa caça a mais derrete a sobra de coelhos,
e os números voltam para perto do equilíbrio.
\end{solution}

\begin{solution}{exo:g5:food-webs:8}
Resposta aberta. Uma teia de quintal defensável: alface e capim
(\omterm{def:g5:food-webs:roles}{produtores}) $\rightarrow$ caracol, pulgão, pardal; caracol
$\rightarrow$ ouriço; pulgão $\rightarrow$ joaninha; pardal e
joaninha $\rightarrow$ gato do bairro ou gavião-miúdo.
\end{solution}

\begin{solution}{exo:g5:food-webs:9}
Primeiro passo: os gafanhotos, menos comidos, aumentam; o gavião,
mais pobre por uma presa, se apoia mais nos coelhos (a seta de
reserva dele). Segundo passo: mais gafanhotos apertam mais o capim;
mais caça a coelhos empurra os coelhos para baixo --- o que alivia o
capim, por sua vez. A teia se dobra e absorve em parte a perda; as
\omterm{def:g5:biodiversity:species}{espécies} sem setas de reserva teriam sofrido mais, e aqui todo mundo
tinha reserva.
\end{solution}

\begin{solution}{exo:g5:food-webs:10}
Gaviões retirados; então os ratos e os ratos-do-campo, livres dos
principais caçadores deles, se multiplicam; então os roedores
multiplicados comem a colheita --- mais grão do que os gaviões jamais
custaram. Quem fez a campanha previu o passo um e parou; a teia
percorreu dois.
\end{solution}

\begin{solution}{exo:g5:food-webs:11}
A perda do carvalho atinge para cima quem depende dele: a \omterm{ex:g2:animals-grow-up:butterfly}{lagarta}
(as \omterm{def:g1:plants-around-us:parts}{folhas} de carvalho eram a única seta dela) passa fome, então o
chapim se apoia inteiramente nos gafanhotos, e então o gavião-miúdo
sente o afinamento --- uma seca de baixo para cima, como o verão
seco. A perda do gavião-miúdo viaja para baixo: os chapins, menos
caçados, aumentam, apertando \omterm{ex:g2:animals-grow-up:butterfly}{lagartas} e gafanhotos --- a história da
retirada do caçador. Aqui a perda do gavião é absorvida com mais
facilidade: as presas dele têm outros caçadores e a perturbação vai
sumindo andar por andar, enquanto quem dependia do carvalho não tinha
seta de reserva nenhuma --- o andar inteiro da \omterm{ex:g2:animals-grow-up:butterfly}{lagarta} se apoiava num
único \omterm{def:g5:food-webs:roles}{produtor}.
\end{solution}
